\section{Mengen }

Georg Cantor beschrieb 1895 eine Menge sinngemäß als Zusammenfassung
bestimmter, wohlunterschiedener Objekte zu einem Ganzen.

\textbf{Beispiele:}

\begin{align*}
 S &:= \{\text{Berlin},\text{ Hamburg},\text{ Köln},\text{ Hannover}\},\\
 \sqrt{2} &\in \mathbb{R},\qquad \{\sqrt{2}\}\subseteq\mathbb{R},\\
 \emptyset &:= \{\}\subseteq\mathbb{R},\\
 \infty &\notin\mathbb{R},\\
 \{1,2,7,\text{Berlin}\} &\nsubseteq\mathbb{R}.
\end{align*}

Die Menge $\mathbb{R}$ umfasst die rationalen und die irrationalen
Zahlen. Beispielsweise gehören $\frac{1}{2}$, $\sqrt{2}$ und $\pi$
zu den reellen Zahlen.

\textbf{Wichtige Zahlenmengen:}

\begin{align*}
 \mathbb{N} &:= \{0,1,2,3,\ldots\},\\
 \mathbb{Z} &:= \{\ldots,-2,-1,0,1,2,\ldots\},\\
 \mathbb{Q} &:= \left\{\frac{p}{q}\mid p\in\mathbb{Z},\;
                       q\in\mathbb{N}\setminus\{0\}\right\},\\
 \mathbb{C} &:= \{a+\mathrm{i}b\mid a,b\in\mathbb{R},\;\mathrm{i}^2=-1\}.
\end{align*}

\textbf{Hinweis:} Die natürlichen Zahlen werden nicht überall gleich
definiert. In diesem Skript gilt $0\in\mathbb{N}$. Alternativ wird
häufig $\mathbb{N}=\{1,2,3,\ldots\}$ und
$\mathbb{N}_0=\{0,1,2,\ldots\}$ verwendet.

Es gelten die Inklusionen
\[
 \mathbb{N}\subseteq\mathbb{Z}\subseteq\mathbb{Q}
 \subseteq\mathbb{R}\subseteq\mathbb{C}.
\]

\textbf{Mathematische Schreibweisen:}

\begin{align*}
 x\in A &\quad\text{bedeutet: $x$ ist Element von $A$,}\\
 x\notin A &\quad\text{bedeutet: $x$ ist kein Element von $A$,}\\
 A\subseteq B &\quad\text{bedeutet: $A$ ist Teilmenge von $B$,}\\
 \emptyset &\quad\text{bezeichnet die leere Menge.}
\end{align*}
In einer Mengenbeschreibung bedeutet $\mid$ so viel wie
„mit der Eigenschaft“.

\subsection{Intervalle}

\textbf{Intervalle:}\\
Seien $a,b\in\mathbb{R}$ mit $a<b$.\\
$\blacktriangleright$ Abgeschlossenes Intervall\\  
\begin{quote}
	$[a,b]:=\{x\in\mathbb{R}\;|\;a\leq x \leq b\}$\vspace*{4mm}\\
	\begin{tikzpicture}
		\node at (2,1) [rotate=0,color=red] {$[$};
		\node at (2,0.5) [rotate=0, color=red] {$a$};
		\draw[->] (0,1) -- (6,1) node [below]{$\mathbb{R}$};
		\draw  (2,1) -- (4,1)[ color=red] ;
		\node at (4,1) [rotate=0,color=red] {$]$};
		\node at (4,0.5) [rotate=0, color=red] {$b$};
	\end{tikzpicture}\\
	z.B.: $[1,3]$\\
	\begin{tikzpicture}
		\node at (2,1) [rotate=0,color=red] {$[$};
		\node at (2,0.5) [rotate=0, color=red] {$1$};
		\draw[->] (0,1) -- (6,1) node [below]{$\mathbb{R}$};
		\draw  (2,1) -- (4,1)[ color=red] ;
		\node at (4,1) [rotate=0,color=red] {$]$};
		\node at (4,0.5) [rotate=0, color=red] {$3$};
	\end{tikzpicture}\\
\end{quote}
$\blacktriangleright$ Rechts halboffenes Intervall\\ 
\begin{quote}
	$[a,b[:=\{x\in\mathbb{R}\;|\;a\leq x < b\}$\vspace*{4mm}\\
	\begin{tikzpicture}
		\node at (2,1) [rotate=0,color=red] {$[$};
		\node at (2,0.5) [rotate=0, color=red] {$a$};
		\draw[->] (0,1) -- (6,1) node [below]{$\mathbb{R}$};
		\draw  (2,1) -- (4,1)[ color=red] ;
		\node at (4,1) [rotate=0,color=black] {$[$};
		\node at (4,0.5) [rotate=0, color=black] {$b$};
	\end{tikzpicture}\\
	$b\notin[a,b[$\\
	(analog: das links halboffene Intervall: $\quad  ]a,b]$) \\
\end{quote}
$\blacktriangleright$ Offenes Intervall\\ 
\begin{quote}
	$]a,b[:=\{x\in\mathbb{R}\;|\;a<x < b\}$\vspace*{4mm}\\
	\begin{tikzpicture}
		\node at (2,1) [rotate=0,color=black] {$]$};
		\node at (2,0.5) [rotate=0, color=black] {$a$};
		\draw[->] (0,1) -- (6,1) node [below]{$\mathbb{R}$};
		\draw  (2,1) -- (4,1)[ color=red] ;
		\node at (4,1) [rotate=0,color=black] {$[$};
		\node at (4,0.5) [rotate=0, color=black] {$b$};
	\end{tikzpicture}\ \\
\end{quote}
$\blacktriangleright$ Spezielle Teilintervalle\\
\begin{quote}
	Geometrisch gedeutet ist dies eine Teilgerade der reellen Zahlen.\vspace*{4mm}\\ $]-\infty,a[:=\{x\in\mathbb{R}\;|\;x<a\}$\vspace*{4mm}\\
	\begin{tikzpicture}
		
		
		\draw[<->] (0,1) -- (6,1) node [below]{$\mathbb{R}$};
		\draw  (0,1) -- (4,1)[ color=red] ;
		\node at (4,1) [rotate=0,color=black] {$[$};
		\node at (4,0.5) [rotate=0, color=black] {$a$};
	\end{tikzpicture}\ \\
\end{quote}
Das offene Intervall $]a,b[$ wird auch als $(a,b)$ geschrieben.\\
Das rechts halboffene Intervall $[a,b[$ wird auch als $[a,b)$ geschrieben.\\
Das links halboffene Intervall $]a,b]$ wird auch als $(a,b]$ geschrieben.\\
\subsection{Durchschnitt, Vereinigung und Differenzmenge}

\textbf{Durchschnitt, Vereinigung und Differenzmenge}\\
$\blacktriangleright$ Durchschnitt\\
\begin{quote}
	
	\tikzset{
		schraffiert/.style={pattern=horizontal lines,pattern color=#1},
		schraffiert/.default=black
	}
	\begin{tikzpicture}
		\def\kreisUeins{(0.5,2) circle [radius=1.5]}
		\def\kreisUzwei{(1.5,2) circle [radius=1.5]}
		\begin{scope} % lokal halten der Wirkung von \clip
			\clip\kreisUeins;
			\draw[schraffiert=green]\kreisUzwei;
		\end{scope}
		\draw\kreisUeins;
		\draw\kreisUzwei;
		\node[anchor=east] at (-0.3,2) {$A$};
		\node[anchor=west] at (2.2,2) {$B$};
	\end{tikzpicture}\\
	
	$\qquad \quad A\cap B$\\
	Beispiel:\\
	Bilde: $[1,2] \cap [2,3]$\\
	\begin{tikzpicture}
		
		
		\draw[->] (0,1) -- (6,1) node [below]{$\mathbb{R}$};
		
		\node at (1,1) [rotate=0,color=black] {$[$};
		\node at (1,0.5) [rotate=0, color=black] {$1$};
		\node at (2,1) [rotate=0,color=black] {$[$};
		\node at (2,0.5) [rotate=0, color=black] {$2$};
		
		
	\end{tikzpicture}\ \\
	\begin{tikzpicture}
		
		
		\draw[->] (0,1) -- (6,1) node [below]{$\mathbb{R}$};
		
		\node at (2,1) [rotate=0,color=black] {$[$};
		\node at (2,0.5) [rotate=0, color=black] {$2$};
		\node at (3,1) [rotate=0,color=black] {$]$};
		\node at (3,0.5) [rotate=0, color=black] {$3$};
		
		
	\end{tikzpicture}\ \\
	$[1,2] \cap [2,3]=\{2\}$\\
\end{quote}
$\blacktriangleright$ Vereinigung\\
\begin{quote}
	
	\tikzset{
		schraffiert/.style={pattern=horizontal lines,pattern color=#1},
		schraffiert/.default=black
	}
	
	\begin{tikzpicture}
		\def\kreisUeins{(0.5,2) circle [radius=1.5]}
		\def\kreisUzwei{(1.5,2) circle [radius=1.5]}
		\begin{scope} % lokal halten der Wirkung von \clip
			%\clip\kreisUeins;
			\draw[schraffiert=green]\kreisUzwei;
			\draw[schraffiert=green]\kreisUeins;
		\end{scope}
		\draw\kreisUeins;
		\draw\kreisUzwei;
		\node[anchor=east] at (-0.3,2) {$A$};
		\node[anchor=west] at (2.2,2) {$B$};
	\end{tikzpicture}\\
	Beispiel:\\
	Bilde: $[1,2]\cup [2,3] $\\
	\begin{tikzpicture}
		
		
		\draw[->] (0,1) -- (6,1) node [below]{$\mathbb{R}$};
		
		\node at (1,1) [rotate=0,color=black] {$[$};
		\node at (1,0.5) [rotate=0, color=black] {$1$};
		\node at (2,1) [rotate=0,color=black] {$]$};
		\node at (2,0.5) [rotate=0, color=black] {$2$};
		
		
	\end{tikzpicture}\ \\
	\begin{tikzpicture}
		
		
		\draw[->] (0,1) -- (6,1) node [below]{$\mathbb{R}$};
		
		\node at (2,1) [rotate=0,color=black] {$[$};
		\node at (2,0.5) [rotate=0, color=black] {$2$};
		\node at (3,1) [rotate=0,color=black] {$]$};
		\node at (3,0.5) [rotate=0, color=black] {$3$};
		
		
	\end{tikzpicture}\ \\
	$[1,2]\cup [2,3]=[1,3]$\\
\end{quote}
\pagebreak
$\blacktriangleright$ Differenzmenge\\
\begin{quote}
	
	\tikzset{
		schraffiert/.style={pattern=horizontal lines,pattern color=#1},
		schraffiert/.default=black
	}
	\begin{tikzpicture}
		\def\kreisUeins{(0.5,2) circle [radius=1.5]}
		\def\kreisUzwei{(1.5,2) circle [radius=1.5]}
		\begin{scope} % lokal halten der Wirkung von \clip
			\clip\kreisUeins;
			\draw[schraffiert=green]\kreisUzwei;
			\draw[schraffiert=green]\kreisUeins;
		\end{scope}
		\draw\kreisUeins;
		\draw\kreisUzwei;
		\fill[white] (1.5,2) circle (1.5);
		\draw (0.5,2) circle (1.5);
		\node[anchor=east] at (-0.3,2) {$A$};
		\node[anchor=west] at (2.2,2) {$B$};
	\end{tikzpicture}\\
	Beispiel:\\
	$[1,3] \; \backslash \;] \frac{3}{2},\frac{5}{2} ]$	\\
	\begin{tikzpicture}
		
		
		\draw[->] (0,1) -- (6,1) node [below]{$\mathbb{R}$};
		
		\node at (1,1) [rotate=0,color=black] {$[$};
		\node at (1,0.5) [rotate=0, color=black] {$1$};
		\node at (3,1) [rotate=0,color=black] {$]$};
		\node at (3,0.5) [rotate=0, color=black] {$3$};
		
		
	\end{tikzpicture}\ \\
	\begin{tikzpicture}
		
		
		\draw[->] (0,1) -- (6,1) node [below]{$\mathbb{R}$};
		
		\node at (1.5,1) [rotate=0,color=black] {$]$};
		\node at (1.5,0.5) [rotate=0, color=black] {$\frac{3}{2}$};
		\node at (2.5,1) [rotate=0,color=black] {$]$};
		\node at (2.5,0.5) [rotate=0, color=black] {$\frac{5}{2}$};
		
		
	\end{tikzpicture}\ \\
	$[1,3] \;\backslash \;] \frac{3}{2},\frac{5}{2} ]=[1,\frac{3}{2}] \cup ]\frac{5}{2}, 3]$	\\
	
	\begin{tikzpicture}
		
		
		\draw[->] (0,1) -- (6,1) node [below]{$\mathbb{R}$};
		
		\node at (1,1) [rotate=0,color=black] {$[$};
		\node at (1,0.5) [rotate=0, color=black] {$1$};
		\node at (1.5,1) [rotate=0,color=black] {$]$};
		\node at (1.5,0.5) [rotate=0, color=black] {$\frac{3}{2}$};
		\node at (2.5,1) [rotate=0,color=black] {$]$};
		\node at (2.5,0.5) [rotate=0, color=black] {$\frac{5}{2}$};
		\node at (3,1) [rotate=0,color=black] {$]$};
		\node at (3,0.5) [rotate=0, color=black] {$3$};
	\end{tikzpicture}\ \\	
\end{quote}
\subsection{Aufgabe zu Mengen 1}

\textbf{Aufgabe:}

Bestimme die folgende Menge $\mathbb{M}$:
\[
 \mathbb{M}=\{x\in\mathbb{Z}\mid 3\text{ teilt }x\}
 \cap\{x\in\mathbb{Z}\mid 4\text{ teilt }x\}.
\]

\textbf{Lösung:}

Die durch 3 bzw. 4 teilbaren ganzen Zahlen bilden die Mengen
\[
 \{3n\mid n\in\mathbb{Z}\}
 \quad\text{bzw.}\quad
 \{4n\mid n\in\mathbb{Z}\}.
\]
Da 3 und 4 teilerfremd sind, beträgt ihr kleinstes gemeinsames
Vielfaches $\operatorname{kgV}(3,4)=12$. Genau die Vielfachen von 12
sind daher durch beide Zahlen teilbar. Folglich gilt:
\[
 \boxed{\mathbb{M}=\{12n\mid n\in\mathbb{Z}\}.}
\]

\subsection{Aufgabe zu Mengen 2}

\textbf{Aufgabe:}

Bestimme die folgende Menge $\mathbb{M}$:

\begin{align*}
	\mathbb{M}
	= \bigcup_{m\in\mathbb{N}\setminus\{0\}}
	\{x\in\mathbb{Q}\mid m\cdot x\in\mathbb{Z}\}
\end{align*}

\textbf{Lösung:}

Betrachten wir zunächst den Fall $m=2$:

\begin{align*}
	M_2=\{x\in\mathbb{Q}\mid 2x\in\mathbb{Z}\}
\end{align*}

Welche rationalen Zahlen ergeben mit 2 multipliziert
eine ganze Zahl?

Beispiele:

\begin{align*}
	\frac{1}{2}\cdot 2 &= 1\\
	-\frac{1}{2}\cdot 2 &= -1\\
	\frac{7}{2}\cdot 2 &= 7
\end{align*}

Damit erhalten wir:

\begin{align*}
	M_2
	= \left\{\frac{n}{2}\mid n\in\mathbb{Z}\right\}
\end{align*}

Für ein beliebiges festes
$m\in\mathbb{N}\setminus\{0\}$ gilt entsprechend:

\begin{align*}
	M_m
	&= \{x\in\mathbb{Q}\mid mx\in\mathbb{Z}\}\\
	&= \left\{\frac{n}{m}\mid n\in\mathbb{Z}\right\}
\end{align*}

Vereinigt man diese Mengen über alle positiven
natürlichen Zahlen $m$, so erhält man sämtliche
rationalen Zahlen.

Denn jede rationale Zahl lässt sich als Bruch
$\frac{n}{m}$ mit $n\in\mathbb{Z}$ und
$m\in\mathbb{N}\setminus\{0\}$ darstellen.

Somit gilt:

\begin{align*}
	\boxed{
		\mathbb{M}
		= \bigcup_{m=1}^{\infty} M_m
		= \bigcup_{m=1}^{\infty}
		\left\{\frac{n}{m}\mid n\in\mathbb{Z}\right\}
		= \mathbb{Q}
	}
\end{align*}



\subsection{Aufgabe zu Mengen 3}

\textbf{Aufgabe:}

Bestimme die folgende Menge:
\[
 \mathbb{M}=\bigcap_{m\in\mathbb{R}}
 \{x\in\mathbb{R}\mid(x-m)(x-2)(x-3)=0\}.
\]

\textbf{Lösung:}

Für $m=1$ gilt:
\[
 \{x\in\mathbb{R}\mid(x-1)(x-2)(x-3)=0\}=\{1,2,3\}.
\]
Für $m=\sqrt{2}$ gilt:
\[
 \{x\in\mathbb{R}\mid(x-\sqrt{2})(x-2)(x-3)=0\}
 =\{\sqrt{2},2,3\}.
\]
Für ein beliebiges $m\in\mathbb{R}$ ist die jeweilige Menge
\[
 M_m=\{x\in\mathbb{R}\mid(x-m)(x-2)(x-3)=0\}
 =\{m,2,3\}.
\]
Dabei gilt auch für $m=2$ oder $m=3$ die Mengenschreibweise
$M_m=\{m,2,3\}$; doppelte Elemente werden in Mengen nicht mehrfach
gezählt.

Die Werte $x=2$ und $x=3$ erfüllen die Gleichung für jedes
$m\in\mathbb{R}$, weil jeweils einer der beiden festen Faktoren
verschwindet. Für ein $x\notin\{2,3\}$ lässt sich dagegen ein
$m\neq x$ wählen. Dann ist keiner der drei Faktoren null.
Somit ist kein weiterer Wert in allen Mengen enthalten.

Es folgt:
\[
 \boxed{\mathbb{M}=\bigcap_{m\in\mathbb{R}}M_m=\{2,3\}.}
\]
